% ----- CHAUPAI 36 -----

\newpage

\subsection*{\deva{षट्त्रिंशत्तम चौपाई} (Thirty-Sixth Chaupai)}
\addcontentsline{toc}{subsection}{\deva{षट्त्रिंशत्तम चौपाई} (Thirty-Sixth Chaupai) :  \deva{संकट कटै...}}

\begin{minipage}[t]{0.55\textwidth}
\vspace{0pt}

{\large \linespread{1.5}\selectfont
\deva{संकट कटै मिटै सब पीरा।}\\
\deva{जो सुमिरै हनुमत बलबीरा॥}
\par}

\vspace{0.5em}

{\large \linespread{1.5}\selectfont
\telu{సంకట కటై మిటై సబ పీరా।}\\
\telu{జో సుమిరై హనుమత బలబీరా॥}
\par}

\vspace{0.5em}

{\large \linespread{1.3}\selectfont
\textit{saṅkaṭa kaṭai miṭai saba pīrā |}\\
\textit{jo sumirai hanumata balabīrā ||}
\par}
\end{minipage}%
\hfill
\begin{minipage}[t]{0.42\textwidth}
\vspace{0pt}
\begin{tcolorbox}[anchorbox]
\textbf{Converting Constraints to Opportunities:} In the *Valmiki Ramayana* (*Sundara Kanda, 48:45-50*), when bound by Indrajit's weapon, Hanuman intentionally chose not to break free immediately. He realized that being bound allowed him to be taken directly to Ravana's court, giving him the perfect opportunity to meet the enemy king face-to-face and deliver a message of peace and warning. It teaches us to convert constraints into opportunities.
\vspace{0.5em}\newline Hanuman strategically accepted his temporary captivity, converting a severe constraint into an opportunity to face the enemy king directly.\newline\null\hfill\href{https://www.valmiki.iitk.ac.in/sloka?field_kanda_tid=5&language=dv&field_sarga_value=48}{\textit{Sundara Kanda, 48:45-50}}
\end{tcolorbox}
\end{minipage}

\vspace{0.5em}
\subsubsection*{\textbf{Visual Rhythm Map:}}

\begin{table}[h!]
\centering
\renewcommand{\arraystretch}{1.2}
{%
\setlength{\tabcolsep}{0pt}%
\begin{tabularx}{\textwidth}{|*{16}{>{\centering\arraybackslash\hsize=1.0\hsize}X|}}
\hline
\multicolumn{4}{|>{\centering\arraybackslash\hsize=4.0\hsize}X|}{\textbf{\guru{सं}\laghu{क}\laghu{ट}} \par (4)} & \multicolumn{3}{>{\centering\arraybackslash\hsize=3.0\hsize}X|}{\textbf{\laghu{क}\guru{टै}} \par (3)} & \multicolumn{3}{>{\centering\arraybackslash\hsize=3.0\hsize}X|}{\textbf{\laghu{मि}\guru{टै}} \par (3)} & \multicolumn{2}{>{\centering\arraybackslash\hsize=2.0\hsize}X|}{\textbf{\laghu{स}\laghu{ब}} \par (2)} & \multicolumn{4}{>{\centering\arraybackslash\hsize=4.0\hsize}X|}{\textbf{\guru{पी}\guru{रा}} \par (4)} \\
\hline
\noalign{\vspace{6pt}}
\hline
\multicolumn{2}{|>{\centering\arraybackslash\hsize=2.0\hsize}X|}{\textbf{\guru{जो}} \par (2)} & \multicolumn{4}{>{\centering\arraybackslash\hsize=4.0\hsize}X|}{\textbf{\laghu{सु}\laghu{मि}\guru{रै}} \par (4)} & \multicolumn{4}{>{\centering\arraybackslash\hsize=4.0\hsize}X|}{\textbf{\laghu{ह}\laghu{नु}\laghu{म}\laghu{त}} \par (4)} & \multicolumn{6}{>{\centering\arraybackslash\hsize=6.0\hsize}X|}{\textbf{\laghu{ब}\laghu{ल}\guru{बी}\guru{रा}} \par (6)} \\
\hline
\end{tabularx}%
}%
\end{table}




\vspace{1em}
\textbf{ \deva{सरल-संस्कृतम्}:}\\
 \deva{तस्य सर्वाणि सङ्कटानि छिद्यन्ते, सर्वाः पीडाः च नश्यन्ति}\\
 \deva{यः बलवीरं (महाबल-शालिनं) हनूमन्तं स्मरति॥}\\


All crises are cut away, and all pain is erased for whoever remembers the mighty and brave Hanuman.


\vspace{1em}
\newpage
\textbf{\deva{प्रतिपदार्थः} (Meaning of each word):}
\begin{table}[h!]
\centering
\renewcommand{\arraystretch}{1.2}
\begin{tabularx}{\textwidth}{@{} l l X X @{}}
\toprule
\textbf{Awadhi} & \textbf{Sanskrit} & \textbf{English} & \textbf{Grammar \& Notes} \\
\midrule
\deva{संकट कटै} / \telu{సంకట కటై} / \textit{saṅkaṭa kaṭai}  &  \deva{सङ्कटं छिद्यते}  &  Crisis is cut  & Present tense `-ai` \\
\deva{मिटै} / \telu{మిటై} / \textit{miṭai}  &  \deva{मृज्यते (नश्यति)}  &  Is erased/mitigated  &   \\
\deva{सब पीरा} / \telu{సబ పీరా} / \textit{saba pīrā}  &  \deva{सर्वाः पीडाः}  &  All pain  & \deva{पीडा $\rightarrow$ पीरा} \\ 
\deva{जो सुमिरै} / \telu{జో సుమిరై} / \textit{jo sumirai}  &  \deva{यः स्मरति}  &  Whoever remembers  & Swara-Bhakti (Smara $\rightarrow$ Sumira) \\ 
\deva{हनुमत} / \telu{హనుమత} / \textit{hanumata}  &  \deva{हनुमन्तम्}  &  Hanuman  &   \\
\deva{बलबीरा} \gotchaicon / \telu{బలబీరా} / \textit{balabīrā}  &  \deva{बलवीरम्}  &  Mighty and Brave  &  \\ 
\bottomrule
\end{tabularx}
\end{table}

\begin{tcolorbox}[poetrybox]
\textbf{The Compound \deva{बलवीर} \gotchaicon:}\\
The Sanskrit original is \deva{बलवीरम्} — a \textit{Dvandva} compound meaning ''one who possesses both \deva{बल} (raw physical strength) and \deva{वीर} (heroic courage).'' 
Tulsidas uses it to make a point: Hanuman is not merely a brute-force warrior (\textit{Bala}), 
nor merely a courageous one (\textit{Vira}), 
but a rare soul in whom both qualities are perfectly fused.\\[0.4em]
The familiar \textbf{\deva{व} $\rightarrow$ \deva{ब} shift} then converts \deva{वीर} into \deva{बीर} in Awadhi, giving us \deva{बलबीरा}. 
The final \deva{-आ} is a standard Awadhi masculine rhyming inflection to complete the 6-matra cadence.
\end{tcolorbox}

\newpage
